\documentclass[10pt,a4paper]{amsart}
\usepackage[margin=19mm]{geometry}
\usepackage{amsmath,amssymb,mathtools}
\usepackage[scr=rsfs,cal=euler]{mathalfa}
\usepackage{xfrac}
\usepackage[unicode,hidelinks,bookmarks=false]{hyperref}
\newcommand{\ie}{i.e.\ }
\newcommand{\cf}{cf.\ }
\newcommand{\resp}{respectively\ }
\newcommand{\Subsection}{\subsection}
\providecommand{\nobreakdash}{\nobreak\hbox{-}}
\setlength{\emergencystretch}{2em}
\multlinegap0pt
\usepackage{bm}
\usepackage{bbm}
\usepackage{dsfont}
\usepackage{xspace}
\usepackage{cancel}
\usepackage{mathtools}
\usepackage{amsthm}
\makeatletter
\@ifundefined{proposition}{\newtheorem{proposition}{Proposition}}{}
\makeatother
\title[The pre-commitment best-choice problem: exact finite-n formu]{The pre-commitment best-choice problem: exact finite-n formulas}
\author{Marcos Costa Santos Carreira}
\date{Source: arXiv:1805.11556v2 (CC BY 4.0)}
\begin{document}
\maketitle

\noindent\emph{Source and licence.} The pre-commitment best-choice problem: exact finite-n formulas. Version arXiv:1805.11556v2 (CC BY 4.0), distributed under the Creative Commons Attribution 4.0 International licence, as recorded in the arXiv licence metadata for this exact version and shown on the abstract page https://arxiv.org/abs/1805.11556v2. Full rights notice: licenses/arxiv-1805.11556-CC-BY-4.0.txt.\par\smallskip

\noindent\textit{Editorial scope.} Counted unit: the Proposition whose source label is \texttt{prop:subopt} in the article (source0.tex lines 380--383, source label \texttt{prop:subopt}), delivered together with the article's own complete original proof of it (source0.tex lines 385--412, one \texttt{proof} environment of 28 lines). No mathematics is written, repaired or re-derived: statement and proof are the article's own text line for line. Required context retained uncounted: none. Every definition, display and label of those lines is retained unchanged. Omitted from this package and not redistributed: 1--379, 384--384, 413--606. Nothing inside a retained excerpt is omitted or altered: every retained body is delivered exactly as the article's lines are written, so no omission receipt of any kind accompanies this package. Citations: the retained excerpts carry no citation command at all, so no bibliography entry of the article is transcribed and no reference is redistributed. Reference adaptation: none was needed and none was made; every reference inside the delivered unit resolves inside it, so no unresolved reference or citation is produced. Printed slips kept literal: no printed word, symbol, hypothesis or number of the counted statement or of its proof was altered. Layout and class: the article's own class is \texttt{article} with packages including inputenc, geometry, amsmath, amssymb, amsthm, graphicx, booktabs, tikz, xcolor, hyperref, authblk, enumitem; the wrapper is the lane's pinned \texttt{amsart} editorial wrapper at 10pt on A4, which restates the theorem-like environments the retained text needs and adds the article's own macro definitions that the retained text needs (none), with extra wrapper packages (bm, bbm, dsfont, xspace, cancel, mathtools). The delivered numbering is the unit's own. Assets: no retained span contains a figure environment, a graphics-inclusion command, a PDF-page inclusion, a table or a TikZ picture; no image file is copied into this package, the package declares no asset dependency and no image file is redistributed with it. Licence: version arXiv:1805.11556v2 of this article is distributed under the Creative Commons Attribution 4.0 International licence, as recorded in the arXiv licence metadata for this exact version and shown on the abstract page https://arxiv.org/abs/1805.11556v2. A full rights notice is delivered at licenses/arxiv-1805.11556-CC-BY-4.0.txt. Characters: every retained excerpt is pure ASCII, so the delivered engine, XeLaTeX with its default Latin Modern text font, reports no missing character.
\begin{proposition}\label{prop:subopt}
For every $n\ge 1$, $v_{pc}(n)\le v(n)$, and for every $n\ge 2$ the inequality is strict.
Consequently $\limsup_{n\to\infty} v_{pc}(n)\le v_\infty=0.580164\ldots$
\end{proposition}

\begin{proof}
\textit{(Subclass bound and a coupling.)} A pre-commitment rule ``stop at the first $r$
with $a_r\ge k_r$'' is a stopping time, so the class is a subset of all stopping rules
and $v_{pc}(n)\le v(n)$. To see where mass is lost, fix a threshold vector $k$ and couple
the pre-commitment rule $\sigma$ with the rule $\tilde\sigma$ that uses the \emph{same}
thresholds but stops only at a draw that is simultaneously above its threshold and a
candidate (running maximum), taking the last draw if no such round occurs. The two agree
until the first crossing $\tau=\min\{r:a_r\ge k_r\}$. If $a_\tau$ is a candidate, both
stop at $\tau$ with identical outcome. If $a_\tau$ is \emph{not} a candidate then
$a_\tau<\max(a_1,\dots,a_{\tau-1})$, so $a_\tau$ is not the overall maximum, $\sigma$
loses with certainty, and $\tilde\sigma$ continues. Hence, with
$F=\{a_\tau\text{ not a candidate}\}$,
\begin{equation}\label{eq:coupling}
  W(\tilde\sigma)-W(\sigma)=\Pr(F)\,\Pr(\tilde\sigma\text{ wins}\mid F)\ \ge\ 0 .
\end{equation}
Taking the supremum over $k$ re-gives $v_{pc}(n)\le v(n)$, and $\limsup_n v_{pc}(n)\le
\lim_n v(n)=v_\infty$.

\textit{(Strictness for finite $n$.)} The optimum $v(n)$ is attained only by rules that
never stop at a non-candidate, since stopping at a non-candidate wins with probability
$0$ and is strictly dominated by continuing (future draws are independent of the past and
carry positive winning probability). Every non-increasing pre-commitment rule with
$k_1<1$ stops at a non-candidate with positive probability: at $n=2$ the region
$\{k_1\le a_1,\ a_1<a_2\}$ has positive area, and for general $n$ the region
$\{k_r\le a_r<\max(a_1,\dots,a_{r-1})\}$ at any round $r\ge 2$ has positive measure
whenever $k_r<1$. (If $k_1=1$, the rule cannot stop early and is suboptimal directly.)
Thus no pre-commitment rule attains $v(n)$ and $v_{pc}(n)<v(n)$ for all $n\ge 2$.
\end{proof}

\end{document}
